\chapter{Motion Estimation}

\section{Table of Contents}

\begin{itemize}
  \item Core Idea
  \item Key Concepts
  \item Theory and Formulas
  \item Visual Schemes
  \item Examples
\end{itemize}

\section{Core Idea}

\textbf{Motion Estimation (ME)} extracts displacement information between video frames. In Multimedia Communications its main role is \textbf{temporal prediction}: if a region in the current frame can be predicted from a reference frame, the codec stores a \textbf{motion vector} and a smaller residual instead of coding the full frame independently.

The course focuses on \textbf{apparent motion}, also called \textbf{optical flow}, not on physical 3D motion. Apparent motion is the 2D displacement of image brightness patterns:

\[
\mathbf{v}(n,m)=
\begin{bmatrix}
u(n,m)\\
v(n,m)
\end{bmatrix}
\]

where $u$ and $v$ are horizontal and vertical pixel displacements.

\begin{calloutbox}{red!70!black}{Physical motion vs. apparent motion}
Physical motion is object movement in 3D space. Apparent motion is movement observed in the image plane. They can differ: a uniformly colored rotating object may have physical motion but no visible optical flow, while a moving light source may create apparent motion without object displacement.
\end{calloutbox}

\begin{center}\includegraphics[width=0.833\linewidth]{ch07/apparent-vs-physical-motion.png}\end{center}
The figure shows why optical flow is an image-domain concept: visible intensity patterns, not real object mechanics, determine the estimated motion.

\section{Key Concepts}

\subsection{Families of Motion Estimation Methods}

\begin{center}\fitwidth{%
\begin{tabular}{|l|l|l|l|}
\hline
\textbf{Family} & \textbf{Output} & \textbf{Model} & \textbf{Typical use} \\ \hline
\textbf{Variational / gradient-based} & dense field, one vector per pixel & smooth optical flow & vision analysis, foundational theory \\ \hline
\textbf{Block matching} & one vector per block & piecewise-constant translation & video compression standards \\ \hline
\textbf{Parametric} & global or regional field & closed-form model, e.g. affine & camera/object motion models \\ \hline
\textbf{Deep learning} & dense learned flow & CNN / recurrent correspondence model & optical-flow benchmarks, analysis tasks \\ \hline
\end{tabular}}
\end{center}

\subsection{Motion-Compensated Prediction}

Given a current frame $f_k$ and a reference frame $f_h$, the motion-compensated prediction is:

\[
\tilde{f}_k(n,m)=
f_h\left(n+u_{h\to k}(n,m),\;m+v_{h\to k}(n,m)\right)
\]

The prediction residual is:

\[
e(n,m)=f_k(n,m)-\tilde{f}_k(n,m)
\]

In compression, good ME should reduce residual energy while keeping the motion-vector field cheap to encode.

\subsection{How to Evaluate a Motion Field}

\begin{center}\fitwidth{%
\begin{tabular}{|l|l|l|}
\hline
\textbf{Criterion} & \textbf{Formula / measure} & \textbf{Meaning} \\ \hline
\textbf{Prediction quality} & $\mathcal{E}=\frac{1}{NM}\sum e^2(n,m)$ & residual MSE after compensation \\ \hline
\textbf{PSNR} & $10\log_{10}\frac{255^2}{\mathcal{E}}$ & quality of the predicted frame \\ \hline
\textbf{MV coding cost} & entropy or coded bits of $(u,v)$ & bit cost of the motion field \\ \hline
\textbf{Complexity} & blocks $\times$ candidates $\times$ cost evaluation & implementation cost \\ \hline
\end{tabular}}
\end{center}

\begin{calloutbox}{red!70!black}{Compression tradeoff}
The best geometric match is not always the best coding decision. A slightly worse prediction can be preferable if its motion vector is much cheaper to transmit.
\end{calloutbox}

\section{Theory and Formulas}

\subsection{Variational Optical Flow}

Gradient-based methods assume \textbf{brightness constancy}:

\[
f_k(n,m)=f_h(n-u,\;m-v)
\]

For small motion, Taylor expansion gives the optical-flow constraint:

\[
\nabla f \cdot \mathbf{v}+f_t=0
\]

or equivalently:

\[
f_xu+f_yv+f_t=0
\]

\begin{calloutbox}{red!70!black}{Aperture problem}
The optical-flow equation gives one scalar equation for two unknowns $(u,v)$. Locally, an edge reveals only the motion component normal to the edge. Additional constraints are needed to estimate a full vector.
\end{calloutbox}

\subsection{Horn and Schunck}

\textit{Horn and Schunck} adds a smoothness prior to the optical-flow equation. It estimates $u$ and $v$ by minimizing:

\[
E(u,v)=
\iint_{\mathcal{R}}
\left[
\left(\nabla f \cdot \mathbf{v}+f_t\right)^2
+\lambda^2
\left(
|\nabla u|^2+|\nabla v|^2
\right)
\right]dn\,dm
\]

The first term enforces brightness constancy. The second term penalizes rapid spatial changes in the velocity field.

\begin{calloutbox}{red!70!black}{Role of $\lambda$}
Small $\lambda$ follows the data more closely and can be noisy. Large $\lambda$ enforces smoother flow but can blur motion boundaries.
\end{calloutbox}

The iterative update used in the slides is:

\[
u^{k+1}=\bar{u}^{k}
-f_x
\frac{f_x\bar{u}^{k}+f_y\bar{v}^{k}+f_t}
{\lambda^2+f_x^2+f_y^2}
\]

\[
v^{k+1}=\bar{v}^{k}
-f_y
\frac{f_x\bar{u}^{k}+f_y\bar{v}^{k}+f_t}
{\lambda^2+f_x^2+f_y^2}
\]

where $\bar{u}^{k}$ and $\bar{v}^{k}$ are local averages.

\begin{center}\includegraphics[width=0.6\linewidth]{ch07/optical-flow-input-frame.png}\end{center}
Input video frame used to illustrate dense optical-flow estimation.

\begin{center}\includegraphics[width=0.6\linewidth]{ch07/horn-schunck-dense-flow.png}\end{center}
Dense optical flow estimated over the frame. The method gives one vector per pixel but tends to smooth across object boundaries.

\subsection{Block Matching}

\textbf{Block matching (BM)} divides the current frame into blocks and assigns one translation vector to each block. For a block $B_{p,q}$ and search window $\mathcal{W}$:

\[
(\hat{i},\hat{j})=
\arg\min_{(i,j)\in\mathcal{W}}
d\left[
\mathbf{f}_k(B_{p,q}),
\mathbf{f}_h(B_{p-i,q-j})
\right]
\]

All pixels inside the block share the same vector:

\[
\forall(n,m)\in B_{p,q}:\quad
\mathbf{v}(n,m)=(\hat{i},\hat{j})
\]

\textbf{Forward motion} uses a previous reference frame ($h=k-1$). \textbf{Backward motion} uses a future reference frame ($h=k+1$).

\begin{center}\includegraphics[width=0.833\linewidth]{ch07/block-matching-reference-current.png}\end{center}
Block matching searches in a reference frame for the block that best predicts the current block.

\subsection{SAD and SSD}

The general norm-based matching cost is:

\[
J_p(i,j)=
\sum_{(n,m)\in B_{p,q}}
\left|
f(n,m,k)-f(n-i,m-j,h)
\right|^p
\]

Two important cases:

\begin{center}\fitwidth{%
\begin{tabular}{|l|l|l|}
\hline
\textbf{Cost} & \textbf{Formula} & \textbf{Properties} \\ \hline
\textbf{SAD} & \$\textbackslash{}sum & f\_k-f\_h \\ \hline
\textbf{SSD} & $\sum (f_k-f_h)^2$ & directly minimizes squared residual energy, but more outlier-sensitive \\ \hline
\end{tabular}}
\end{center}

\begin{center}\includegraphics[width=0.833\linewidth]{ch07/sad-vs-ssd-motion-vectors.png}\end{center}
SAD can give a more regular motion-vector field, while SSD may slightly improve PSNR because it matches the MSE objective.

\subsection{Regularized Motion Estimation}

Motion vectors can be regularized by penalizing vectors that differ from neighboring vectors:

\[
J_{\text{REG}}(i,j)=
\left\|
\mathbf{f}_k(B_{p,q})-\mathbf{f}_h(B_{p-i,q-j})
\right\|_p^p
+\lambda R(i,j)
\]

Modern codecs use a rate-distortion form:

\[
J(v)=D(v)+\lambda_{ME}R(v)
\]

where $D(v)$ is prediction distortion and $R(v)$ is the number of bits needed to encode the vector.

\begin{center}\includegraphics[width=0.833\linewidth]{ch07/regularized-motion-vectors.png}\end{center}
Regularization reduces MVF entropy by preferring smoother vectors, at the cost of a small prediction-quality loss.

\subsection{Block Size Tradeoff}

\begin{center}\fitwidth{%
\begin{tabular}{|l|l|l|l|}
\hline
\textbf{Block size} & \textbf{Complexity} & \textbf{MV coding cost} & \textbf{Prediction error} \\ \hline
larger blocks & lower & lower & higher \\ \hline
smaller blocks & higher & higher & lower \\ \hline
\end{tabular}}
\end{center}

Large blocks are efficient but cannot model detailed object boundaries. Small blocks are flexible but produce more vectors and higher overhead.

\subsection{Search Strategies}

For a full search window:

\[
\mathcal{W}=\{-A,\ldots,A\}\times\{-B,\ldots,B\}
\]

If $A=B=7$, full search tests:

\[
(2A+1)^2=15^2=225
\]

candidates per block.

\begin{center}\fitwidth{%
\begin{tabular}{|l|l|l|l|}
\hline
\textbf{Strategy} & \textbf{Main idea} & \textbf{Typical tests in source example} & \textbf{Risk} \\ \hline
\textbf{Full Search} & test every candidate & 225 & high cost \\ \hline
\textbf{Three Step Search} & decreasing step size & about 25 & local minima \\ \hline
\textbf{Diamond Search} & large diamond, then small diamond & about 23 & local minima \\ \hline
\textbf{Hexagon Search} & large hexagon, then small hexagon & about 17 & local minima \\ \hline
\textbf{TZSearch} & predictors plus adaptive search & variable & more complex control \\ \hline
\end{tabular}}
\end{center}

\begin{center}\includegraphics[width=0.833\linewidth]{ch07/diamond-search-pattern.png}\end{center}
Diamond Search uses a large pattern to move across the cost surface and a small pattern for final refinement.

\begin{center}\includegraphics[width=0.833\linewidth]{ch07/hexagon-search-pattern.png}\end{center}
Hexagon Search reduces new tests per iteration and is used in H.264/AVC-style fast ME.

\subsection{Subpixel Precision}

Motion is not limited to integer pixels. Codecs usually refine integer vectors at half-pixel and quarter-pixel precision.

For bilinear interpolation:

\[
f(n+a,m+b)=
(1-a)(1-b)x
+a(1-b)y
+(1-a)bz
+abw
\]

where $x,y,z,w$ are the four neighboring integer samples.

\begin{center}\includegraphics[width=0.833\linewidth]{ch07/subpixel-motion-grid.png}\end{center}
Subpixel vectors improve prediction quality but require interpolation filters and extra search steps.

\subsection{Variable Block Size}

Variable-size BM splits a block only when the rate-distortion cost improves:

\begin{enumerate}
  \item Estimate the best vector for the current block and compute $J(v)=D+\lambda R$.
  \item Split the block into four sub-blocks.
  \item Estimate each sub-block and compute $J_{\text{sub}}=\sum_i J_i$.
  \item Keep the split only if $J_{\text{sub}}<J(v)$.
  \item Recurse until no gain or the minimum block size is reached.
\end{enumerate}

This is the practical solution used by modern standards: flat regions keep large blocks, while object boundaries use smaller partitions.

\subsection{Parametric Motion Models}

Parametric ME represents the entire MVF with a small number of parameters.

\textbf{Translation:}

\[
\mathbf{v}(\mathbf{p})=
\begin{bmatrix}
b_1\\
b_2
\end{bmatrix}
\]

\textbf{Affine model:}

\[
\mathbf{v}(\mathbf{p})=
\mathbf{b}+\mathbf{B}\mathbf{p}
=
\begin{bmatrix}
b_1\\
b_2
\end{bmatrix}
+
\begin{bmatrix}
b_3 & b_4\\
b_5 & b_6
\end{bmatrix}
\mathbf{p}
\]

The affine model has six parameters and can represent translation, zoom, rotation, shear, and their combinations.

\begin{center}\includegraphics[width=0.6\linewidth]{ch07/affine-motion-field-translation.png}\end{center}
Pure translation gives a constant vector field.

\begin{center}\includegraphics[width=0.6\linewidth]{ch07/affine-motion-field-zoom.png}\end{center}
Zoom and related affine effects are represented by a position-dependent field.

Parameter estimation can be:

\begin{itemize}
  \item \textbf{indirect}: estimate a dense flow first, then fit parameters by least squares;
  \item \textbf{direct}: estimate parameters by minimizing the optical-flow residual or a SAD/SSD prediction error directly.
\end{itemize}

For an affine model, direct optical-flow fitting uses:

\[
\pi^*=
\arg\min_\pi
\sum_{(n,m)\in\mathcal{R}}
\left[
u_\pi(n,m)f_x+v_\pi(n,m)f_y+f_t
\right]^2
\]

\subsection{Deep Learning Methods}

Deep learning treats optical flow as a learned correspondence problem. Training can be supervised with synthetic ground truth or unsupervised with photometric reconstruction losses.

\begin{center}\fitwidth{%
\begin{tabular}{|l|l|l|l|}
\hline
\textbf{Model} & \textbf{Year} & \textbf{Main contribution} & \textbf{Limitation} \\ \hline
\textbf{FlowNet} & 2015 & first end-to-end CNN for dense optical flow & weak on fine details and small displacements \\ \hline
\textbf{FlowNet 2.0} & 2017 & stacked networks and small-displacement module & large memory footprint \\ \hline
\textbf{PWC-Net} & 2018 & pyramids, warping, cost volume & coarse-to-fine detail loss \\ \hline
\textbf{RAFT} & 2020 & all-pairs correlation and recurrent GRU updates & high memory cost \\ \hline
\end{tabular}}
\end{center}

\begin{center}\includegraphics[width=0.833\linewidth]{ch07/flownet-architecture.png}\end{center}
FlowNet introduced end-to-end CNN optical flow, including variants with simple concatenation and explicit feature correlation.

\begin{center}\includegraphics[width=0.833\linewidth]{ch07/flownet2-stacked-architecture.png}\end{center}
FlowNet 2.0 improves accuracy by stacking networks and specializing modules for different displacement regimes.

\begin{calloutbox}{red!70!black}{Compression vs. analysis}
Learned dense flow dominates many analysis benchmarks. Block matching remains central in video compression because it is simple, controllable, standardizable, and naturally tied to rate-distortion optimization.
\end{calloutbox}

\section{Visual Schemes}

\subsection{Motion-Compensated Coding Loop}

\begin{center}\footnotesize\itshape Diagram (Mermaid source):\end{center}
\begin{verbatim}
flowchart LR
    A[Current frame f_k] --> B[Block partition]
    B --> C[Motion estimation]
    R[Reference frame f_h] --> C
    C --> D[Motion vectors]
    D --> E[Motion compensation]
    R --> E
    E --> F[Prediction f_hat_k]
    A --> G[Residual e]
    F --> G
    D --> H[MV coding]
    G --> I[Residual coding]
    H --> J[Bitstream]
    I --> J
\end{verbatim}

The loop shows why ME is a codec tool: it creates a predictor, then only vectors and residual information must be transmitted.

\subsection{Method Overview}

\begin{center}\includegraphics[width=\linewidth]{ch07/motion-estimation-methods-overview.png}\end{center}
The overview connects the main ME families: optical flow, block matching, parametric models, and learned methods.

\subsection{Variable Block-Size Decision}

\begin{center}\footnotesize\itshape Diagram (Mermaid source):\end{center}
\begin{verbatim}
flowchart TD
    A[Block b] --> B[Estimate vector for b]
    B --> C[Cost J = D + lambda R]
    C --> D[Split into 4 sub-blocks]
    D --> E[Estimate sub-block vectors]
    E --> F[Cost J_sub = sum J_i]
    F --> G{J_sub < J?}
    G -- yes --> H[Keep split and recurse]
    G -- no --> I[Keep parent block]
\end{verbatim}

This decision balances lower residual distortion against the extra rate required by more vectors and partition syntax.

\section{Examples}

\subsection{SAD vs. SSD on Flower Garden}

\begin{center}\fitwidth{%
\begin{tabular}{|l|l|l|l|}
\hline
\textbf{Criterion} & \textbf{MV rate} & \textbf{Prediction PSNR} & \textbf{Interpretation} \\ \hline
SSD & 2143 bits & 22.46 dB & best MSE-oriented prediction \\ \hline
SAD & 2103 bits & 22.30 dB & slightly lower PSNR, lower MV cost \\ \hline
\end{tabular}}
\end{center}

\begin{calloutbox}{green!55!black}{Exam takeaway}
SSD optimizes squared error directly, so it can improve PSNR. SAD is often preferred in practice because it is cheaper, more robust to outliers, and can produce a more regular vector field.
\end{calloutbox}

\subsection{Regularized SSD}

\begin{center}\fitwidth{%
\begin{tabular}{|l|l|l|l|}
\hline
\textbf{Method} & \textbf{MV rate} & \textbf{Prediction PSNR} & \textbf{Effect} \\ \hline
SSD & 2143 bits & 22.46 dB & no MV regularization \\ \hline
Regularized SSD & 2008 bits & 22.35 dB & lower MV rate with small PSNR loss \\ \hline
\end{tabular}}
\end{center}

Rate reduction:

\[
\frac{2143-2008}{2143}\approx 6.3\%
\]

The example shows the principle behind rate-distortion ME: a small distortion increase can be accepted if it saves enough motion-vector bits.

\subsection{Search Complexity Example}

For $A=B=7$, full search tests $225$ candidates. Fast search methods in the source example need about:

\begin{center}\fitwidth{%
\begin{tabular}{|l|l|l|}
\hline
\textbf{Method} & \textbf{Tests} & \textbf{Approximate saving} \\ \hline
2D-log / Three Step & 25 & 89\% fewer \\ \hline
Diamond Search & 23 & 90\% fewer \\ \hline
Hexagon Search & 17 & 92\% fewer \\ \hline
\end{tabular}}
\end{center}

Fast methods are efficient because natural-video cost surfaces are often locally well behaved, but they do not guarantee the global minimum.

\subsection{Final Comparison}

\begin{center}\fitwidth{%
\begin{tabular}{|l|l|l|}
\hline
\textbf{Method} & \textbf{Strength} & \textbf{Weakness} \\ \hline
Horn-Schunck & dense, convex, elegant theory & over-smooths boundaries \\ \hline
Full-search BM & global optimum in search window & expensive \\ \hline
Fast BM & very low complexity & possible local minima \\ \hline
Variable-size BM & good codec tradeoff & partition overhead and control complexity \\ \hline
Affine ME & compact global/regional field & needs suitable region/model \\ \hline
FlowNet / PWC-Net / RAFT & accurate dense flow for analysis & training data, memory, generalization \\ \hline
\end{tabular}}
\end{center}

\begin{calloutbox}{red!70!black}{What to remember}
Motion estimation is a rate-distortion problem in video coding, not only a geometric matching problem. The best codec decision balances residual energy, motion-vector rate, search complexity, block partitioning, and interpolation precision.
\end{calloutbox}
